\documentclass[10pt,a4paper]{amsart}
\usepackage[margin=19mm]{geometry}
\usepackage{amsmath,amssymb,mathtools}
\usepackage[scr=rsfs,cal=euler]{mathalfa}
\usepackage{xfrac}
\usepackage[unicode,hidelinks,bookmarks=false]{hyperref}
\newcommand{\ie}{i.e.\ }
\newcommand{\cf}{cf.\ }
\newcommand{\resp}{respectively\ }
\newcommand{\Subsection}{\subsection}
\providecommand{\nobreakdash}{\nobreak\hbox{-}}
\setlength{\emergencystretch}{2em}
\multlinegap0pt
\usepackage{mathrsfs}
\usepackage{epsfig}
\usepackage{latexsym}
\usepackage{bbm}
\usepackage{caption2}
\usepackage{color}
\newcommand{\thetaop}{\operatorname{\theta}}
\newcommand {\cross } {\nu_D }
\newcommand {\function } {f_D}
\usepackage{txfonts}
\usepackage{mathrsfs}
\newtheorem{theorem}{Theorem}[section]
\newtheorem{lemma}[theorem]{Lemma}
\newtheorem{cor}[theorem]{Corollary}
\newtheorem{prop}[theorem]{Proposition}
\newtheorem{conj}[theorem]{Conjecture}
\newtheorem{remark}[theorem]{Remark}
\newtheorem{definition}[theorem]{Definition}
\newtheorem{exam}[theorem]{Example}
\newtheorem{problem}[theorem]{Problem}
\newtheorem{question}[theorem]{Question}
\usepackage{latexsym}
\newcommand{\QED}{\hfill $\Box$}
\newcommand{\Za}{\hbox{\ bf Z}_a}
\newcommand{\N}{\hbox{\bf N}}
\newcommand{\Z}{\hbox{\bf Z}}
\newcommand{\subgp}[1]{\langle{#1}\rangle}
\newcommand{\F}{{\mathbb F}}
\newcommand{\qfd}{\hfill $\Box$}
\begin{document}
\title[The universal zero-sum invariant and weighted zero-sum for i]{The universal zero-sum invariant and weighted zero-sum for infinite abelian groups II}
\author{Guoqing Wang}
\date{}
\maketitle
\noindent\emph{Source and licence.} The universal zero-sum invariant and weighted zero-sum for infinite abelian groups II. Version arXiv:2607.02132v1 (CC BY 4.0). This material is distributed under the Creative Commons Attribution 4.0 International licence, http://creativecommons.org/licenses/by/4.0/, as recorded in the arXiv OAI licence metadata for this exact version and shown on the abstract page https://arxiv.org/abs/2607.02132v1. Full rights notice: licenses/arxiv-2607.02132-CC-BY-4.0.txt.\par\smallskip
\noindent\textit{Editorial scope.} Counted unit: the original \texttt{theorem} environment \texttt{thm:kernel-cover} at source lines 464--494 together with its complete proof at lines 496--544; the delivered document keeps source lines 423--545 so that every referenced label and every citation resolves inside it. Native editable source is the paper's own concatenated TeX; the AMS wrapper is editorial formatting only. No publisher class, upstream script or unrelated artwork is redistributed.
\section{Finiteness of Weighted Davenport constant}






We begin this section with some terminologies. Let $\{F_i:i\in I\}$ be a family of abelian groups, and let $G$ be an
abelian group. Let
$\psi_i:F_i\longrightarrow G \qquad (i\in I)$
be homomorphisms. By the universal property of the direct sum, there
exists a unique homomorphism
$
  \psi:\bigoplus\limits_{i\in I}F_i\longrightarrow G
$
such that
\begin{equation}\label{equation universal map}
  \psi\circ\iota_i=\psi_i
  \qquad\text{for every }i\in I,
\end{equation}
where
$
  \iota_i:F_i\longrightarrow\bigoplus\limits_{t\in I}F_t
$
is the canonical injection. We call $\psi$ the \emph{universal
homomorphism induced by the family $\{\psi_i:i\in I\}$}, denoted
$$\psi=\operatorname{ind}[\{\psi_i:i\in I\}].$$
Now let $n\geq 1$ and let $F_1=\cdots=F_n=F$. We identify
$\bigoplus\limits_{i=1}^n F_i$ with $F^n$. For any family of homomorphism
$\{\psi_i\}_{i\in [1,n]}$, also denoted as $(\psi_1,\ldots,\psi_n)$, we can write
$\operatorname{ind}[\{\psi_i:i\in [1,n]\}]$ as $\operatorname{ind}[(\psi_1,\ldots,\psi_n)].$
Therefore, for any $\boldsymbol{\eta}=(\psi_1,\ldots,\psi_n)\in \operatorname{Hom}(F,G)^n$ and any $(x_1,\ldots,x_n)\in F^n$,
we have
\begin{equation}\label{eq:induced-explicit}
\operatorname{ind}[\boldsymbol{\eta}]\left((x_1,\ldots,x_n)\right)=\operatorname{ind}[\boldsymbol{\eta}]\left(\sum\limits_{i=1}^n \iota_i(x_i)\right)=\sum\limits_{i=1}^n  \operatorname{ind}[\boldsymbol{\eta}]\left(\iota_i(x_i)\right)=\sum\limits_{i=1}^n \psi_i(x_i).
\end{equation}
Note that if $(\psi_1,\ldots,\psi_n)\neq (\mathbf 0,\ldots,\mathbf 0)$, then the induced universal
homomorphism $\operatorname{ind}[(\psi_1,\ldots,\psi_n)]$ is a nonzero homomorphism of $F^n$ to $G$.

Then we shall prove our main theorem for this section.


\begin{theorem}\label{thm:kernel-cover}
Let $F$ and $G$ be abelian groups, let
$\varnothing\neq\Psi\subseteq\operatorname{Hom}(F,G)\setminus\{\mathbf 0\},$
and let $n\ge 1$. Then the following two statements are equivalent.
\begin{enumerate}
  \item[\rm (i)] ${\rm D}_{\Psi}(G)\le n$.

  \item[\rm (ii)] The group $F^n$ is covered by the kernels of the
  nonzero universal homomorphisms induced by the tuples in
  $(\Psi\cup\{\mathbf 0\})^n$; more precisely,
  $$F^n=\bigcup_{\boldsymbol{\eta}\in
      (\Psi\cup\{\mathbf 0\})^n\setminus\{(\mathbf 0,\ldots,\mathbf 0)\}}
    \ker\bigl(\operatorname{ind}[\boldsymbol{\eta}]\bigr).$$
\end{enumerate}

Consequently,
\begin{equation}\label{eq:D-minimum}
 {\rm D}_{\Psi}(G)
 =
 \min\left\{
 n\in\mathbb N:
 F^n
 =
 \bigcup_{\boldsymbol{\eta}\in
   (\Psi\cup\{\mathbf 0\})^n\setminus
   \{(\mathbf 0,\ldots,\mathbf 0)\}}
 \ker\bigl(\operatorname{ind}[\boldsymbol{\eta}]\bigr)
 \right\},
\end{equation}
where the minimum of the empty set is understood to be $\infty$.
\end{theorem}

\begin{proof}
Assume first that ${\rm D}_{\Psi}(G)\le n$. We take an arbitrary element
$(x_1,\ldots,x_n)\in F^n.$
Consider the sequence
$T=x_1\cdot\ldots\cdot x_n$
over $F$. By the definition of ${\rm D}_{\Psi}(G)$, the sequence $T$
has a nonempty $\Psi$-weighted zero-sum subsequence. Hence there exist
a {\bf nonempty} set $I\subseteq[1,n]$ and homomorphisms
$\psi_i\in\Psi \ (i\in I)$
such that
$\sum_{i\in I}\psi_i(x_i)=0_G.$ For each $j\in[1,n]\setminus I$, we set $\psi_j=\mathbf 0$. Denote
$\boldsymbol{\eta}=(\psi_1,\ldots,\psi_n).$
Since $I\neq\varnothing$, the tuple $\boldsymbol{\eta}$ is not the
zero tuple, i.e., $\boldsymbol{\eta}\in
   (\Psi\cup\{\mathbf 0\})^n\setminus
   \{(\mathbf 0,\ldots,\mathbf 0)\}$. By \eqref{eq:induced-explicit}, we derive that
$\operatorname{ind}[\boldsymbol{\eta}]\left((x_1,\ldots,x_n)\right)=\sum_{i=1}^n\psi_i(x_i)=\sum_{i\in I}\psi_i(x_i)=0_G.$ That is,
$(x_1,\ldots,x_n)\in
  \ker\bigl(\operatorname{ind}[\boldsymbol{\eta}]\bigr).$
By the arbitrariness of choosing $(x_1,\ldots,x_n)$ from $F^n$, Statement \emph{(ii)} holds. This proves that Statement (i) implies Statement (ii).

Now we show that Statement (ii) implies Statement (i). Suppose that (ii) holds.  Let
$T=x_1\cdot\ldots\cdot x_n$
be an arbitrary sequence over $F$ of length $n$.  Since $(x_1,\ldots,x_n)\in F^n$, it follows that there
exists
$\boldsymbol{\eta}=(\psi_1,\ldots,\psi_n)
  \in
  (\Psi\cup\{\mathbf 0\})^n
  \setminus\{(\mathbf 0,\ldots,\mathbf 0)\}$
such that
$(x_1,\ldots,x_n)
  \in
  \ker\bigl(\operatorname{ind}[\boldsymbol{\eta}]\bigr).$
Set
$J=\{j\in[1,n]:\psi_j\neq\mathbf 0\}.$
Since $\boldsymbol{\eta}$ is not the zero tuple, we have $J\neq
\varnothing$. Note that $\psi_j\in\Psi$ for each $j\in J$, and that
$0_G=\operatorname{ind}[\boldsymbol{\eta}]
      \left((x_1,\ldots,x_n)\right)=
      \sum_{i=1}^n\psi_i(x_i) =
      \sum_{j\in J}\psi_j(x_j).$
Thus
$\prod\limits_{j\in J}x_j$
is a nonempty $\Psi$-weighted zero-sum subsequence of $T$. By the arbitrariness of choosing $T\in \mathcal{F}(F)$ with $|T|=n$, it follows that
${\rm D}_{\Psi}(G)\le n.$

The formula \eqref{eq:D-minimum} follows immediately from the
equivalence of (i) and (ii).
\end{proof}


\end{document}
